\subsection{One-endpoint estimation} \label{subsec:One-endpoint_estimation}

Fix a one-endpoint coordinate $U^{(a)}_{\ell,j,\alpha}$ with label $\alpha=(s_\alpha,i_\alpha)$. It is \textit{high-degree} when $D_a\ge\kappa/2$ and \textit{low-degree} otherwise. We handle these cases in Sections~\ref{sec:highdegree} and~\ref{Sec:lowdegree}, respectively.

We test the length and literal sign as each clause arrives and use quantum queries to test its prefix counts. A successful query selects one clause; the sub-sketch then counts its exact suffix degrees and uses them with the candidate prefix values to test the bucket. Its output is scaled by $M$ or $M/2$ as specified below.

\paragraph{Prefix counts.}
Recall that $B_a>D_{a+1}+1$. For nonnegative integers $d,c$, define
\[
  \Gamma_a(d,c):=d+B_a c.
\]
In a high-degree sketch, $c$ is the sampled positive-prefix count $s_{v,a}^{\le C}$; in a low-degree sketch, it is the exact positive-prefix degree $p_v^{\le C}$.

For every $q\in\{0,\ldots,\kappa-1\}$, integer $c\ge0$, and degree $D_a<d\le D_{a+1}$,
\begin{equation}\label{eq:one-endpoint-ordinary-high-id}
\1[\Gamma_a(d,c)\ge B_a q+D_a+1]
-\1[\Gamma_a(d,c)\ge B_a q+D_{a+1}+1]
=
\1[c=q].
\end{equation}
Likewise, for integers $D_a<d,d_0\le D_{a+1}$ and $c,c_0\ge0$,
\begin{equation}\label{eq:one-endpoint-ordinary-low-id}
\1[\Gamma_a(d,c)\ge B_a c_0+d_0]
-\1[\Gamma_a(d,c)\ge B_a c_0+d_0+1]
=
\1[d=d_0,\ c=c_0].
\end{equation}
For every integer $q\ge1$, the same bound on $B_a$ gives
\begin{equation}\label{eq:one-endpoint-cumulative-high-id}
\1[\Gamma_a(d,c)\ge qB_a]=\1[c\ge q]
\qquad(D_a<d\le D_{a+1},\ c\ge0).
\end{equation}

\paragraph{Impersonation.}
Set
\[
H_a:=
\begin{cases}
\kappa, & D_a\ge\kappa/2,\\
D_{a+1}, & D_a<\kappa/2.
\end{cases}
\]
We say that an endpoint with prefix degree $d$ \emph{impersonates degree interval $a$} if $d>D_{a+1}$ and
\begin{equation}\label{eq:scale-a-impersonation-def}
d\in
\bigcup_{h=1}^{H_a}
\bigl[hB_a+D_a+1,\ hB_a+D_{a+1}\bigr].
\end{equation}
Let $\mathsf{Imp}^{a}_{C,r}$ denote this event for endpoint $(C,r)$. The condition includes every nonzero threshold difference from outside the queried degree interval, although an endpoint satisfying it may give a zero difference.

\begin{lemma}[Impersonation]\label{lem:low-degree-singleton-impersonation}
Assume $D_a<\kappa/2$. Fix a queried low-degree level $B_a c_0 +d_0$, where $D_a<d_0\le D_{a+1}$ and $0\le c_0 \le d_0$. Let $d^\ast,c^\ast$ be nonnegative integers. If
\[
\Gamma_a(d^\ast,c^\ast)=B_a c_0 +d_0
\]
and $d^\ast\notin(D_a,D_{a+1}]$, then $d^\ast>D_{a+1}$ and the endpoint with prefix degree $d^\ast$ impersonates degree interval $a$.
\end{lemma}

\begin{proof}
If $d^\ast\le D_a$, equality modulo $B_a$ gives $d^\ast=d_0>D_a$, a contradiction. Thus $d^\ast>D_{a+1}$. Writing $h:=c_0-c^\ast$, we have
\[
d^\ast=hB_a+d_0>D_{a+1},\qquad
1\le h\le c_0\le d_0\le D_{a+1}=H_a.
\]
Since $D_a<d_0\le D_{a+1}$, this gives
$d^\ast\in[hB_a+D_a+1,hB_a+D_{a+1}]$, as required.
\end{proof}

\subsubsection{High-degree case} \label{sec:highdegree}

For each $q\in\{0,1,\ldots,\kappa\}$, we maintain an independent
\textit{$q$-Sketch}. The case $q=\kappa$ accounts for all endpoints with
$s_{v,a}^{\le C}\ge\kappa$ through the capped value
$\min\{s_{v,a}^{\le C},\kappa\}$, so high-degree one-endpoint coordinates need no
separate overflow term.

Each $q$-Sketch consists of two independent sub-sketches, called the \textit{Upper-$q$-Sketch} and the \textit{Lower-$q$-Sketch}. Their outputs are combined to give an estimate for the high-degree coordinate.

The quantum register uses the following basis labels. We use the global register size $M=\lceil \gamma_\eps m\rceil$, with $\gamma_\eps$ chosen large enough for the fixed $\eps>0$.
\begin{enumerate}
    \item The fixed anchor basis state $\ket{\bot}$;
    \item Spare labels: $S_1,\ldots, S_M$;
    \item $H$-labels: $H(v,t)$, where $v\in [n]$, $t\in [M]$;
    \item $E$-labels: $E(v,t)$, where $v\in [n]$, $t\in [M]$;
    \item $H$-garbage-labels: $G_{H}(v,t)$,where $v\in [n]$, $t\in [M]$;
    \item $E$-garbage-labels: $G_{E}(v,t)$,where $v\in [n]$, $t\in [M]$;
\end{enumerate}
Besides, we maintain a pointer $\mathcal{P}$ for the spare labels in a classical register, which costs $O(\log M)=O(\log m)$ bits. Because the basis state is specified by a variable index and a level in $[M]$, this quantum register costs $O(\log (nm))=O(\log n )$ qubits, for $m=\mathrm{poly} (n)$.

\paragraph{What the labels are for.}
The spare labels make each update a permutation and are not queried. Each literal copy on $u$ applies one shift to the $H(u,\cdot)$ labels, which are used to test prefix total degree.

The $E$-labels are used to test thresholds of
\[
\Gamma_a(d_u^{\le C},s_{u,a}^{\le C})
=d_u^{\le C}+B_a s_{u,a}^{\le C}.
\]
Each literal copy applies one shift, and each selected positive copy applies $B_a$ additional shifts. The garbage labels make these updates permutations on a finite set of basis labels.

\begin{definition}[Increment operation]
Given a variable $u\in[n]$ and the current  pointer $\mathcal P$, define $\inc{H,u,1}$ as follows.
\begin{enumerate}
    \item Apply the cyclic permutation of basis labels whose cycle is
    \[
    S_{\mathcal P}\to H(u,1)\to H(u,2)\to \cdots \to H(u,M)
    \to G_H(u,M)\to \cdots \to G_H(u,1)\to S_{\mathcal P}.
    \]
    We call this one elementary shift. Its quantum part is unitary because it is
    a permutation of basis labels.
    \item Move the pointer to the next spare label. If $\mathcal P<M$,
set $\mathcal P\leftarrow \mathcal P+1$. If $\mathcal P=M$, set
$\mathcal P\leftarrow 1$. Thus the pointer runs through
$S_1,S_2,\ldots,S_M$ and then starts again from $S_1$.
\end{enumerate}
For any integer $r\ge1$, $\inc{H,u,r}$ means repeating this elementary shift $r$ times, updating $\mathcal P$ after each repetition, that is, \[
  \mathrm{inc}(H,u,r):=\underbrace{\mathrm{inc}(H,u,1)\circ\cdots\circ\mathrm{inc}(H,u,1)}_{r\text{ times}}.
\]
For $E$-labels, we use the same shift operation on the chain $E(u,\cdot)$; in particular, $\inc{E,u,B_a}$ means applying this one-step shift $B_a$ times, so it advances $\mathcal P$ by $B_a$.
\end{definition}

We describe the Lower-$q$-Sketch for $q\ge1$ first, then give the changes for the other sub-sketches. Set $\mathsf{Tgt}^{\ell,j,\alpha}_{C}:=0$ if $|C|\ne\ell$; otherwise set
\[
\mathsf{Tgt}^{\ell,j,\alpha}_{C}:=\1[s_j(C)=s_\alpha].
\]


%\subsubsection*{Algorithm Description}


\medskip

Now we are ready to describe our algorithm.
\paragraph{Initialization.} The quantum register is initialized as: 
\[
  \Qstate{T_0}
  =
  \frac{|\perp\rangle+|S_1\rangle+\cdots+|S_{M-1}\rangle}{\sqrt M},
\]
where $T_0:=\{S_1,S_2,\dots,S_{M-1}\}$ is the initial live set.
Initialize $\mathcal{P}=1$.

When a clause $C$ arrives, the quantum algorithm performs two types of operations in turn: update and query.

\paragraph{Update.} Decompose the clause $C$ into literal copies and handle them one by one. For each literal copy $(C,r,h)$, set $u:=v_r(C)$ and let $T$ be the live set. Perform the following increment operations in turn. Each operation applies a unitary permutation to the quantum register and updates the pointer as specified above.

\medskip
\textbf{Step-1.} $
      \mathrm{inc}(H,u,1).
    $

\medskip    
\textbf{Step-2.}      $
      \mathrm{inc}(E,u,1).
    $
    
\medskip    
\textbf{Step-3.} If this literal copy $(C,r,h)$ is positive and $f_a(C,r,h)=1$,  perform   $
      \mathrm{inc}(E,u,B_a).
    $ 

\paragraph{Query.} If $\mathsf{Tgt}^{\ell,j,\alpha}_{C}=0$, do nothing and end. Otherwise, write $v:=v_j(C)$ for its variable index. Perform the pair-query $$(H(v,D_{a}+1),  E(v,qB_a ))$$ on the current quantum register. Denote the output of pair-query by $\chi_{C}^{\mathrm L}\in \{-1,0,+1\}$.    

\medskip
\textbf{Branch 1:} if $\chi_{C}^{\mathrm L} = 0$, we say this query \textit{fails}. Then do nothing and end.

\medskip
\textbf{Branch 2:} if  $\chi_{C}^{\mathrm L}\in \{-1,+1\}$, we say this query is \textit{successful}. Then this quantum register performs no further updates or queries. During the rest of the stream, the sub-sketch only keeps the following classical data: \begin{enumerate}
    \item store the variable $v$ and clause $C$;
    \item store the nonzero measurement outcome $\chi_C^{\mathrm L}$;
    \item store $q$;
    \item in the following streaming of clauses, count the suffix total degree  $d_{v}^{>C}$ and suffix positive count $p_{v}^{>C}$ accurately.
\end{enumerate}
The classical part costs $O(\log (n+m))$ bits. We say a sub-sketch is \textit{active} if the quantum part has not ended, equivalently, if it has not entered Branch 2 yet.

\paragraph{Output.} If a high-degree sub-sketch succeeds on $C$, then after the rest of the stream has been processed it has the exact suffix counts $p_v^{>C}$ and $d_v^{>C}$. For every $r\in\{0,\ldots,\kappa\}$, define
\[
F_r(C)
:=
\1 \left[
\clip
\left(
\frac{2\left(\frac{2D_a}{\kappa}r+p_v^{>C}\right)}
     {D_a+d_v^{>C}}
-1+g(v)
\right)
\in I_{i_\alpha}
\right].
\]
Finally,  the Lower-$q$-Sketch $(q\ge 1)$ outputs 
\[
  Y_q^{\mathrm L}(C):=\frac M2\cdot \chi_{C}^{\mathrm L}\cdot (F_q(C)-F_{q-1}(C)). 
\]
In all other cases,  the Lower-$q$-Sketch $(q\ge 1)$ outputs $0$.

\medskip

For the other sub-sketches, only the following operations and outputs change.

\paragraph{Upper-$q$-Sketch ($q\ge 1$).} In the query step, we perform the pair-query on \[
(H(v,D_{a+1}+1), 
E(v,qB_a ))
\]
rather than $(H(v,D_{a}+1), E(v,qB_a ))$ for Lower-$q$-Sketch.

Similarly, define $\chi^{\mathrm U}_{C}\in \{-1,0,1\}$ for the output of the pair-query. In its successful branch, the Upper-$q$-Sketch stores the nonzero value of $\chi_C^{\mathrm U}$ together with the same classical data. For Branch $2$, define the final output as
\[
  Y_q^{\mathrm U}(C):=\frac M2\cdot \chi_{C}^{\mathrm U}\cdot (F_q(C)-F_{q-1}(C)).
\]
In all other cases, it outputs  $0$.

\paragraph{Lower-$0$-Sketch and Upper-$0$-Sketch.} These sub-sketches use the anchor, spare labels, $H$-labels, and $H$-garbage-labels, and apply only $\inc{H,u,1}$.

In the query step, the Lower-$0$-Sketch performs a single-query on $H(v,D_a+1)$ and the Upper-$0$-Sketch performs a single-query on $
H(v,D_{a+1}+1)$.

Define $\xi_{C}^{\mathrm L}$ and $\xi_{C}^{\mathrm U}$ to denote the outputs of the single-query from Lower-$0$-Sketch and  Upper-$0$-Sketch  respectively. In a successful branch the sketch stores this nonzero output, which equals $1$, and the same classical data and suffix counters as above. Define the final outputs as \[
  Y_0^{\mathrm L}(C):=M\cdot \xi_{C}^{\mathrm L}\cdot F_0(C),
  \qquad
  Y_0^{\mathrm U}(C):=M\cdot \xi_{C}^{\mathrm U}\cdot F_0(C),
\]
In all other cases, output $0$.







\medskip
Next, we prove the correctness of the algorithm and show that its output provides a good estimate of the coordinate $U^{(a)}_{\ell,j,\alpha}$.


\paragraph{Correctness.}
Fix a sub-sketch and a label $F(u,t)$ that it queries, where $F$ is $H$ or $E$. For this fixed family and variable, the queried level $t$ is fixed, and at least one ordinary $+1$ update occurs between consecutive queries of $F(u,t)$. Immediately before a query, let $d_u$ count the ordinary literal-copy updates applied to this chain so far and $c_u$ count its additional $B_a$-shift updates. The elementary shifts inside a $B_a$-shift update are not counted again in $d_u$. Set $c_u=0$ for a family with no additional updates.

\begin{samepage}
\begin{lemma}\label{lem:nonwrapping-threshold-invariant}
Suppose $1\le t<M$ and, before each query of $F(u,t)$, the sub-sketch has used fewer than $M$ elementary shifts in total. Then, immediately before each such query,
\[
F(u,t)\in T
\quad\Longleftrightarrow\quad
d_u+B_a c_u\ge t.
\]
\end{lemma}

The proof of Lemma~\ref{lem:nonwrapping-threshold-invariant} is in Appendix~\ref{app:one-endpoint-threshold-proof}.
\end{samepage}

\begin{remark}
Lemma~\ref{lem:nonwrapping-threshold-invariant} also applies to the $Z_1$- and $Z_2$-labels in Section~\ref{subsec:Two-endpoint_estimation}: each chain has a fixed queried level, and each queried clause updates that chain before the query.
\end{remark}

The next lemma bounds both the queried levels and the number of shifts, including the low-degree sketches described below.

\begin{lemma}[No-wrap event for one-endpoint sketches]\label{lem:good_label_budget}
For each high-degree interval $a$, let $X_a$ count all positive literal copies selected by $f_a$ over the whole stream. For a sufficiently large constant $K_\eps$, define
\[
\mathcal G_1:=
\bigcap_{a:\,D_a\ge\kappa/2}
\left\{X_a\le K_\eps\frac{m}{D_a}\right\}.
\]
Then $\Pr[\mathcal G_1]\ge1-1/160$. Choosing $\gamma_\eps$ in $M=\lceil\gamma_\eps m\rceil$ sufficiently large ensures that every queried level satisfies $1\le t<M$ and, on $\mathcal G_1$, every one-endpoint sub-sketch performs fewer than $M$ elementary shifts over the whole stream.
\end{lemma}

The proof of Lemma~\ref{lem:good_label_budget} is in Appendix~\ref{app:one-endpoint-budget-proof}.

The updates remain permutations even if $\mathcal G_1$ fails. On $\mathcal G_1$, Lemma~\ref{lem:nonwrapping-threshold-invariant} identifies every queried label with its prefix-count threshold.

\paragraph{Estimation.}
Fix public randomness $\omega=(\{f_a\},g,\{B_a\})$ satisfying $\mathcal G_1$. All expectations below average over the quantum measurements conditional on $\omega$. The lower and upper sketches select the degree interval; a telescoping sum over $q$ then tests the bucket.

Extend every sub-sketch output by $0$ to clauses with $\mathsf{Tgt}^{\ell,j,\alpha}_{C}=0$. For each clause $C$, define
\begin{equation}\label{eq:estimation}
\widehat U^{(a)}_{\ell,j,\alpha}(C)
:=
Y_0^{\mathrm L}(C)-Y_0^{\mathrm U}(C)
+\sum_{q=1}^{\kappa}\bigl(Y_q^{\mathrm L}(C)-Y_q^{\mathrm U}(C)\bigr),
\end{equation}
and set
\[
\widehat U^{(a)}_{\ell,j,\alpha}
:=\sum_{C\in\Phi_{\le\ell_0}}\widehat U^{(a)}_{\ell,j,\alpha}(C).
\]
Fix a clause $C$ with $\mathsf{Tgt}^{\ell,j,\alpha}_{C}=1$ and write $v:=v_j(C)$. When $a=a(C,j)$, the high-degree prefix approximations give
\[
\widetilde d_v^{\le C}=D_a,\qquad
\widetilde p_v^{\le C}
=\frac{2D_a}{\kappa}\min\{s_{v,a}^{\le C},\kappa\}.
\]
Using $F_q(C)$ defined in the output step, we obtain
\begin{equation}\label{eq:Fq}
\begin{aligned}
\1[\widetilde\alpha_j(C)=\alpha]
&=F_{\min\{s_{v,a}^{\le C},\kappa\}}(C)\\
&=F_0(C)+\sum_{q=1}^{\kappa}
\bigl(F_q(C)-F_{q-1}(C)\bigr)\1[s_{v,a}^{\le C}\ge q].
\end{aligned}
\end{equation}

The following probability cancels the normalization of the active quantum register.

\begin{lemma}[Active probability]\label{lem:active_prob}
Fix a sub-sketch and public randomness $\omega$. Conditioned on this sub-sketch being active just before its query on $C$, let $T$ be its live set and $M'=1+|T|$. Then
\[
\Pr[\text{this sub-sketch is active just before its query on }C\mid\omega]
=\frac{M'}{M}.
\]
\end{lemma}

\begin{proof}
Once $\omega$ is fixed, the updates and previous failed queries determine the live set on the active branch. List the earlier queries as $\text{Query}_1,\ldots,\text{Query}_r$, and let $M_i$ be one plus the number of live labels just before $\text{Query}_i$ on this branch. Set $M_{r+1}=M'$; initially $M_1=M$.

A failed query deletes zero, one, or two live labels. Updates preserve their number, so the conditional probability of remaining active is
\[
\Pr[\text{active after }\text{Query}_i\mid\omega,\ \text{active before }\text{Query}_i]
=1-\frac{M_i-M_{i+1}}{M_i}
=\frac{M_{i+1}}{M_i}.
\]
Multiplying these probabilities gives $M'/M$.
\end{proof}

\paragraph{For $q\in\{1,\ldots,\kappa\}$.}
Conditioned on the Lower-$q$-Sketch being active before $C$, let $T$ be its live set and $M'=1+|T|$. By \eqref{eq:pair-query-mean}, its pair-query on $(H(v,D_a+1),E(v,qB_a))$ gives
\[
\mathbb E[\chi_C^{\mathrm L}\mid\omega,\ \text{Lower-$q$-Sketch active before }C]
=\frac{2}{M'}\1[H(v,D_a+1)\in T]\1[E(v,qB_a)\in T].
\]
Lemma~\ref{lem:nonwrapping-threshold-invariant} replaces the live-label indicators by prefix-count inequalities. Multiplying the query mean by $\frac M2(F_q(C)-F_{q-1}(C))$ and the active probability $M'/M$ from Lemma~\ref{lem:active_prob} gives
\[
\begin{aligned}
\mathbb E[Y_q^{\mathrm L}(C)\mid\omega]
={}&\bigl(F_q(C)-F_{q-1}(C)\bigr)
\1[\mathsf{Tgt}^{\ell,j,\alpha}_{C}=1,\ d_v^{\le C}\ge D_a+1]\\
&\cdot\1[\Gamma_a(d_v^{\le C},s_{v,a}^{\le C})\ge qB_a].
\end{aligned}
\]
Applying the same argument separately to the Upper-$q$-Sketch gives
\[
\begin{aligned}
\mathbb E[Y_q^{\mathrm U}(C)\mid\omega]
={}&\bigl(F_q(C)-F_{q-1}(C)\bigr)
\1[\mathsf{Tgt}^{\ell,j,\alpha}_{C}=1,\ d_v^{\le C}\ge D_{a+1}+1]\\
&\cdot\1[\Gamma_a(d_v^{\le C},s_{v,a}^{\le C})\ge qB_a].
\end{aligned}
\]
Their difference is zero outside $(D_a,D_{a+1}]$. Inside this interval, \eqref{eq:one-endpoint-cumulative-high-id} gives $\Gamma_a(d_v^{\le C},s_{v,a}^{\le C})\ge qB_a\iff s_{v,a}^{\le C}\ge q$, so
\begin{equation}\label{eq:bigq}
\begin{aligned}
\mathbb E[Y_q^{\mathrm L}(C)-Y_q^{\mathrm U}(C)\mid\omega]
={}&\bigl(F_q(C)-F_{q-1}(C)\bigr)\\
&\cdot\1[\mathsf{Tgt}^{\ell,j,\alpha}_{C}=1,\ D_a<d_v^{\le C}\le D_{a+1},\ s_{v,a}^{\le C}\ge q].
\end{aligned}
\end{equation}
Thus high-degree one-endpoint estimation has no impersonation error.

\paragraph{For $q=0$.}
For the Lower-$0$-Sketch with live set $T$ and $M'=1+|T|$, \eqref{eq:single-query-mean} gives
\[
\mathbb E[\xi_C^{\mathrm L}\mid\omega,\ \text{Lower-$0$-Sketch active before }C]
=\frac{\1[H(v,D_a+1)\in T]}{M'}.
\]
Multiplying by $MF_0(C)$ and its active probability $M'/M$ from Lemma~\ref{lem:active_prob}, and applying the same argument separately to the Upper-$0$-Sketch, gives
\begin{equation}\label{0q}
\mathbb E[Y_0^{\mathrm L}(C)-Y_0^{\mathrm U}(C)\mid\omega]
=F_0(C)\1[\mathsf{Tgt}^{\ell,j,\alpha}_{C}=1,\ D_a<d_v^{\le C}\le D_{a+1}].
\end{equation}
Summing the conditional expectations over $q$ and using the telescoping identity~\eqref{eq:Fq} gives
\begin{equation}\label{eq:one-endpoint-high-coordinate}
\mathbb E[\widehat U^{(a)}_{\ell,j,\alpha}(C)\mid\omega]
=\1[\mathsf{Tgt}^{\ell,j,\alpha}_{C}=1,\ D_a<d_v^{\le C}\le D_{a+1},\ \widetilde\alpha_j(C)=\alpha].
\end{equation}
Summing over clauses yields
\[
\mathbb E[\widehat U^{(a)}_{\ell,j,\alpha}\mid\omega]
=U^{(a)}_{\ell,j,\alpha}.
\]

\subsubsection{Low-degree case} \label{Sec:lowdegree}

For a low-degree coordinate $U^{(a)}_{\ell,j,\alpha}$, there are $O_\eps(1)$ candidate prefix-degree pairs $(d,c)$, with integers $D_a<d\le D_{a+1}$ and $0\le c\le d$. For each pair, we maintain a $(d,c)$-Sketch.

Each $(d,c)$-Sketch consists of independent Lower-$(d,c)$-Sketch and Upper-$(d,c)$-Sketch sub-sketches. Their quantum registers have the same basis labels. The quantum register and update operations are the same as in the high-degree case, except for the positive-count update described next.

For the update step, Step-3 is replaced as follows: if this literal copy is positive, perform
$
      \mathrm{inc}(E,u,B_a).
    $

For the query step, the Lower-$(d,c)$-Sketch performs a single-query on \[
E(v, B_a c+d )
\] and the Upper-$(d,c)$-Sketch performs single-query on \[
E(v, B_a c+d+1 ).
\] 

The difference of their expected outputs tests $\Gamma_a(d_v^{\le C},p_v^{\le C})=B_ac+d$. Since it does not separately test the prefix degree interval, an endpoint with degree outside $(D_a,D_{a+1}]$ may contribute. Lemma~\ref{lem:low-degree-singleton-impersonation} includes every such contribution in $\mathsf{Imp}^{a}_{C,j}$.

\paragraph{Output.}
After a successful query on $C$, the corresponding Lower-$(d,c)$-Sketch or Upper-$(d,c)$-Sketch stores the nonzero single-query output, which equals $1$, and computes
\[
F_{d,c}(C)
:=
\1\!\left[
\clip\!\left(
\frac{2(c+p_v^{>C})}{d+d_v^{>C}}-1+g(v)
\right)\in I_{i_\alpha}
\right]
\]
using the candidate prefix values $(d,c)$ in Definition~\ref{def:weighted-pseudobias}. Define
\[
Y_{d,c}^{\mathrm L}(C):=M\cdot \xi^{\mathrm L}_{d,c}(C)\cdot F_{d,c}(C),
\qquad
Y_{d,c}^{\mathrm U}(C):=M\cdot \xi^{\mathrm U}_{d,c}(C)\cdot F_{d,c}(C),
\]
where $\xi^{\mathrm L}_{d,c}(C),\xi^{\mathrm U}_{d,c}(C)\in\{0,1\}$ are the single-query outputs of the Lower-$(d,c)$-Sketch and Upper-$(d,c)$-Sketch, respectively. Extend both outputs by $0$ to clauses with $\mathsf{Tgt}^{\ell,j,\alpha}_{C}=0$.

\paragraph{Estimation and Analysis.}
Fix $\omega$ satisfying $\mathcal G_1$. For each clause $C$, define
\[
\widehat U^{(a)}_{\ell,j,\alpha}(C)
:=\sum_{d=D_a+1}^{D_{a+1}}\sum_{c=0}^{d}
\bigl(Y_{d,c}^{\mathrm L}(C)-Y_{d,c}^{\mathrm U}(C)\bigr),
\]
and set
\[
\widehat U^{(a)}_{\ell,j,\alpha}
:=\sum_{C\in\Phi_{\le\ell_0}}\widehat U^{(a)}_{\ell,j,\alpha}(C).
\]
Fix a clause $C$ with $\mathsf{Tgt}^{\ell,j,\alpha}_{C}=1$ and write $v:=v_j(C)$. Conditioned on the Lower-$(d,c)$-Sketch being active before $C$, let its live set be $T$ and $M'=1+|T|$. Equation~\eqref{eq:single-query-mean} gives the single-query mean; replacing its live-label indicator by the prefix-count inequality from Lemma~\ref{lem:nonwrapping-threshold-invariant} gives
\[
\begin{aligned}
&\mathbb E[Y_{d,c}^{\mathrm L}(C)\mid\omega,\ \text{Lower-$(d,c)$-Sketch active before }C]\\
&\qquad=\frac{M}{M'}F_{d,c}(C)
\1[\Gamma_a(d_v^{\le C},p_v^{\le C})\ge B_ac+d].
\end{aligned}
\]
Multiplying by its active probability $M'/M$, and applying the same argument separately to the Upper-$(d,c)$-Sketch, yields
\[
\mathbb E[Y_{d,c}^{\mathrm L}(C)-Y_{d,c}^{\mathrm U}(C)\mid\omega]
=F_{d,c}(C)\1[\Gamma_a(d_v^{\le C},p_v^{\le C})=B_ac+d].
\]
If $\mathsf{Imp}^{a}_{C,j}$ does not occur, Lemma~\ref{lem:low-degree-singleton-impersonation} excludes equality outside $(D_a,D_{a+1}]$, while \eqref{eq:one-endpoint-ordinary-low-id} identifies the exact prefix-degree pair inside this interval. Hence
\[
\Gamma_a(d_v^{\le C},p_v^{\le C})=B_ac+d
\quad\Longleftrightarrow\quad
d_v^{\le C}=d,\quad p_v^{\le C}=c.
\]
For this pair, $(d,c)$ equals the prefix approximations, so it gives the same pseudo-bias; with sign $s_\alpha$, $F_{d,c}(C)=\1[\widetilde\alpha_j(C)=\alpha]$ by Definitions~\ref{def:weighted-pseudobias} and~\ref{def:weighted-pseudosnapshot}. Therefore
\[
\mathbb E[\widehat U^{(a)}_{\ell,j,\alpha}(C)\mid\omega]
=\1[D_a<d_v^{\le C}\le D_{a+1},\ \widetilde\alpha_j(C)=\alpha]
\]
whenever $\mathsf{Imp}^{a}_{C,j}$ does not occur.

For any fixed clause, at most one candidate pair satisfies $\Gamma_a(d_v^{\le C},p_v^{\le C})=B_ac+d$: equality for two pairs implies $d-d'=B_a(c'-c)$, and $|d-d'|<B_a$ forces $c=c'$ and $d=d'$. Thus the expected contribution of all $(d,c)$-sketches lies in $[0,1]$. An impersonating endpoint has $d_v^{\le C}>D_{a+1}$ and is not counted by the target coordinate. Clauses with $\mathsf{Tgt}^{\ell,j,\alpha}_{C}=0$ have zero output, so
\begin{equation}\label{eq:one-endpoint-low-coordinate}
\begin{aligned}
\left|
\mathbb E[\widehat U^{(a)}_{\ell,j,\alpha}(C)\mid\omega]
-\1\!\left[\substack{\mathsf{Tgt}^{\ell,j,\alpha}_{C}=1,\\
D_a<d_v^{\le C}\le D_{a+1},\
\widetilde\alpha_j(C)=\alpha}\right]
\right|
\le\1[\mathsf{Tgt}^{\ell,j,\alpha}_{C}=1,\ \mathsf{Imp}^{a}_{C,j}].
\end{aligned}
\end{equation}
Summing over clauses gives
\begin{equation}\label{eq:one-endpoint-low-coordinate-total}
\left|
\mathbb E[\widehat U^{(a)}_{\ell,j,\alpha}\mid\omega]
-U^{(a)}_{\ell,j,\alpha}
\right|
\le\sum_{C\in\Phi_{\le\ell_0}}
\1[\mathsf{Tgt}^{\ell,j,\alpha}_{C}=1,\ \mathsf{Imp}^{a}_{C,j}].
\end{equation}
